\documentclass[11pt,a4paper,oneside]{ltjsbook}

\usepackage[haranoaji]{luatexja-preset}
\usepackage{amsmath,amssymb,mathtools}
\usepackage{booktabs,longtable,array}
\usepackage{geometry}
\usepackage{xcolor}
\usepackage{listings}
\usepackage{tikz}
\usetikzlibrary{arrows.meta,positioning,fit,shapes.geometric,calc}
\usepackage[most]{tcolorbox}
\usepackage{hyperref}
\usepackage{url}
\usepackage{enumitem}
\usepackage{fancyhdr}

\geometry{top=24mm,bottom=28mm,left=24mm,right=24mm,headheight=15pt,footskip=12mm}
\hypersetup{
  colorlinks=true,
  linkcolor=blue!45!black,
  urlcolor=blue!55!black,
  pdftitle={type-pm-mech2 入門},
  pdfauthor={type-pm-mech2 project}
}

\definecolor{CodeBg}{RGB}{246,248,250}
\definecolor{DeepBlue}{RGB}{33,72,117}
\definecolor{DeepGreen}{RGB}{35,105,82}
\definecolor{Warm}{RGB}{150,79,28}
\definecolor{SoftRed}{RGB}{148,55,55}

\lstset{
  basicstyle=\ttfamily\small,
  backgroundcolor=\color{CodeBg},
  frame=single,
  framesep=5pt,
  rulecolor=\color{black!18},
  columns=fullflexible,
  keepspaces=true,
  showstringspaces=false,
  breaklines=true,
  xleftmargin=4pt,
  xrightmargin=4pt
}

\newtcolorbox{pointbox}[1]{
  enhanced,
  colback=blue!3,
  colframe=DeepBlue,
  title={#1},
  fonttitle=\bfseries,
  boxrule=0.7pt,
  arc=2mm
}
\newtcolorbox{warningbox}[1]{
  enhanced,
  colback=orange!4,
  colframe=Warm,
  title={#1},
  fonttitle=\bfseries,
  boxrule=0.7pt,
  arc=2mm
}
\newtcolorbox{definitionbox}[1]{
  enhanced,
  colback=green!3,
  colframe=DeepGreen,
  title={#1},
  fonttitle=\bfseries,
  boxrule=0.7pt,
  arc=2mm
}

\newcommand{\code}[1]{\texttt{#1}}
\newcommand{\checks}{\mathrel{\Leftarrow}}
\newcommand{\TyInt}{\mathsf{Int}}
\newcommand{\Matcher}{\mathsf{Matcher}}
\newcommand{\Slot}{\mathsf{Slot}}
\newcommand{\Prod}{\mathsf{Prod}}
\newcommand{\fn}{\mathbin{\to}}

\pagestyle{fancy}
\fancyhf{}
\fancyhead[L]{\small type-pm-mech2 入門}
\fancyhead[R]{\small\leftmark}
\fancyfoot[C]{\thepage}
\renewcommand{\headrulewidth}{0.3pt}

\title{\Huge type-pm-mech2 入門\\[4mm]
  \Large 制約生成，\code{let}の完全性，matcherまで}
\author{type-pm-mech2 project}
\date{2026年8月20日版\\
  基準実装：\code{main}（本文作成開始時のcommit \code{7fe112f}）}

\begin{document}

\frontmatter
\maketitle

\chapter*{この文書について}
\addcontentsline{toc}{chapter}{この文書について}

この文書は，Lean 4リポジトリ\code{type-pm-mech2}の現行実装を，
型推論や定理証明に初めて触れる読者にも追えるように説明する．
中心となる問いは次の三つである．

\begin{enumerate}[leftmargin=2em]
  \item 式を調べると得られる，仮の結果型\code{target}，必ず満たす等式\code{hard}，
        後で分類する使用可能性の要求\code{pending}とは何か．
  \item 一つの定義を複数の型で使える多相\code{let}で右辺を先に閉じると，なぜ
        「型付けが存在すれば推論が必ず成功する」
        という性質（完全性）の証明が難しくなるのか．
  \item 現行実装はその問題をどう解き，どこまでを証明済みとしているのか．
\end{enumerate}

後半では，Paper 1用の検証例（fixture）として実装されている，値の分解方法を記述する
再帰的multiset matcher，
とくに\code{multisetWithListArgument = lambda list. multiset}という補助式の意味と，
この式が「制約からまだ決まらない型の形を先回りして推測しない」方針
（no-guess方針）の下で，単独では推論に失敗する理由も説明する．

\begin{warningbox}{旧リポジトリと隣接する論文について}
この文書は\code{type-pm-mech2}のLeanコードを正本とする．
兄弟リポジトリ\code{type-pm-mech}は旧実装であり，旧\code{SourceTyping}，
\code{inferRaw}，\code{wBridgeCheck}，terminal audit（推論後の再検査）を
現行体系へ持ち込まない．また，\code{../type-pm-paper}のTeXには旧実装を説明する
記述が残っているため，現行実装のAPI一覧としては使わない．
現状確認には本書，設計と進捗を統合したリポジトリ直下の\code{README.md}，
およびLeanコードを用いる．ここでAPIとは，ほかのファイルから
利用するために公開された定義や定理の入口をいう．
\end{warningbox}

\tableofcontents

\mainmatter

\chapter{まず全体像をつかむ}

\section{何を機械化しているのか}

\code{type-pm-mech2}は，Type-PMという型体系をLean 4で機械化するプロジェクトである．
ここで\emph{機械化}とは，型や推論規則をプログラムとして定義するだけでなく，
「推論成功なら規則上も型が付く」ことや「規則上型が付くなら推論が成功する」ことまで，
Leanの定理として検査することをいう．前者を健全性，後者を完全性と呼び，後で改めて定義する．

通常の関数型言語では，型推論，つまりプログラムを調べてその型を自動的に求める処理は，
式に\code{Int}や$\alpha\fn\beta$のような型を付ける．$\alpha$と$\beta$はまだ分からない型を
表す型変数であり，矢印は「$\alpha$型の値を受け取り$\beta$型の値を返す関数」を表す．
Type-PMではさらにmatcherを扱う．
matcherは，値をどのように分解・観察できるかを記述する実行時の値である．
そのため型にも次の二つの見方が現れる．

\begin{description}[leftmargin=8em,style=nextline]
  \item[matcher producer]
    実際に作られたmatcher値の型．本書では概略$\Matcher\,\kappa\,\tau$と書く．
  \item[matcher slot]
    matcherを使う場所が要求する型．本書では概略$\Slot\,\kappa\,\tau$と書く．
\end{description}

$\tau$は対象値の型であり，$\kappa$はcapabilityである．
capabilityは，matcherが入力値のどの形を観察できるかを表す型情報である．
producerからslotへの変換は単なる型等式ではないため，普通の単一化だけでなく
checkingという別の処理が必要になる．これが後で\code{pending}を分離する理由である．

\section{先に覚える基本語}

この先で繰り返し使う語を，コード上の英語名より先に日本語で説明する．

\begin{description}[leftmargin=10em,style=nextline]
  \item[source式]
    型を調べる対象となる入力言語の式である．本書で単に「式」と書くときも，通常はこの式を指す．
  \item[context]
    式のその場所から参照できる変数と，各変数に与えられた型を並べた環境である．
    たとえば無名関数（lambda）の本体では，lambdaが束縛した引数の型がcontextへ追加される．
    自由変数とは，その式自身のlambdaや\code{let}では導入されず，外側のcontextから参照する変数である．
  \item[代入]
    型変数を具体的な型で一斉に置き換える対応である．たとえば
    $[\alpha\mapsto\TyInt]$は，$\alpha$を\code{Int}で置き換える．
  \item[単一化]
    複数の型等式を同時に成り立たせる代入を計算する処理である．
  \item[関数適用]
    関数を引数へ使う式である．たとえば$f\;x$は関数$f$を引数$x$へ適用する．
  \item[制約block]
    一つの式についてまとめて解く制約の一まとまりである．現行実装では，生の結果型
    \code{target}，必ず満たす等式\code{hard}，後で分類する使用可能性の要求\code{pending}を持つ．
    三つの詳しい意味は次章で説明する．
  \item[型scheme]
    型変数を利用時ごとに選べる形で束ね，一つの定義を複数の型で使えるようにした多相型である．
    この束ねる操作を量化といい，コード名は\code{Scheme}である．
  \item[signature]
    \code{nil}や\code{cons}のようにdata値を作る名前（data constructor）や，\code{add}のような
    組込み演算（primitive）と，それらの型schemeを登録した有限の表である．
  \item[公開推論]
    利用者が式の型を求めるために呼ぶ，リポジトリの正式な推論関数\code{infer}である．
  \item[supply]
    次に作る新しい型変数とcapability変数へ何番を与えるかを記録する二つのカウンタである．
  \item[well-formedなsupply]
    contextにすでに現れるすべての自由変数より後の番号から，新しい変数名を割り当てるsupplyである．
    これにより，生成した新しい名前がcontext中の古い名前と衝突しない．
  \item[closure]
    一つの制約blockを解いて得た代入と結果型をまとめた記録である．
    実行時に関数本体と環境を組にした「関数closure」とは別物である．
  \item[回帰検証]
    具体例と期待結果を固定し，将来の変更で以前の正しい動作が壊れていないかを確かめる検査である．
\end{description}

証明済みの性質を表す四つの語も，次の意味で用いる．

\begin{description}[leftmargin=8em,style=nextline]
  \item[健全性]
    公開推論が成功したなら，推論器とは独立に定義した型付け規則でも型が付くこと．
  \item[完全性]
    推論器とは独立な型付けが存在するなら，公開推論が必ず成功すること．
  \item[主要性]
    公開推論が返す型が最も一般的で，ほかの型付け結果をその型への代入で得られること．
  \item[受理同値]
    二つの制約blockについて，一方を解けることと他方を解けることが同値であること．
\end{description}

\section{M0からM5まで}

実装は段階に分けられている．2026年8月20日時点の要点は表\ref{tab:milestones}のとおりである．

\begin{longtable}{>{\bfseries}p{9mm}p{42mm}p{18mm}p{70mm}}
\caption{実装段階の要約}\label{tab:milestones}\\
\toprule
段階 & 主題 & 状態 & 意味 \\
\midrule
\endfirsthead
\toprule
段階 & 主題 & 状態 & 意味 \\
\midrule
\endhead
M0 & 型，代入，局所checking & 完了 & matcherを含む基礎的な型合成とcheckingを検証済み．\\
M1 & lambda，application，制約block & 完了 & 独立した型付け，停止する単一化，推論の健全性・完全性・主要性を検証済み．\\
M2 & 型scheme，多相\code{let} & 完了 & 一般の入れ子\code{letE}を含むM2--M3のsource式で完全性・受理同値・主要性を検証済み．\\
M3 & data，primitive，signature & 一部 & Bool/List，値を作るconstructor，組込み演算primitive，\code{ifE}，有限signatureを実装済み．M4と結ぶ掲載例全体の回帰が残る．\\
M4 & pattern，matcher，\code{fix} & 一部 & 最終構文と統合推論に加え，pattern function本体の照合を外側から隔離する実行時nodeであるMNodeとPaper 2の\code{pair}回帰を実装済み．source宣言とruntime本体表の双方向対応も\code{pair}で証明した．各runtime本体は保存schemeの標準的な一つの具体化について検査証拠を持つが，全具体化に対する主要性までは未証明である．生成変数の由来を追う証明は，公開M4 elaborationまで完了した．大きい検証例の正確な公開推論とM4全域の完全性・主要性が残る．\\
M5 & 評価と型安全性 & 一部 & 順序付き探索，型付き環境，通常・再帰closure，任意の型付き関数位置を持つapplication，\code{map}まで実装済み．ここで型安全性とは，型の付いた実行が型の不整合で停止しない性質である．source/runtime定義表の整合性証明を必須にするMNode全式評価器と，Paper 2の全式回帰も実装済みだが，評価器全体の関係的意味論と型安全性は未証明である．source多相\code{let}から実行時型付けへの橋とuser-defined matcherの型安全性も残る．\\
\bottomrule
\end{longtable}

\begin{pointbox}{「M2が完了」の正確な意味}
M2では，独立な型付けがあれば公開推論が必ず成功すること（完全性）と，公開推論の型が
ほかのすべての型付け結果の元になること（主要性）を証明済みである．ただし推論を始める
カウンタ（開始supply）は，その場所で参照できる変数の環境（context）にすでに現れる
すべての自由変数より後から，新しい番号を割り当てなければならない．この安全な開始条件を
「contextに対してwell-formedである」と呼ぶ．公開推論はこの条件を満たす
\code{context.initialSupply}から必ず始める．条件を破る任意の開始supplyは未解決の宿題ではなく，
APIの仕様が受け付ける範囲の外である．
\end{pointbox}

\section{三つの層を混同しない}

現行実装には，役割の異なる三つの層がある．

\begin{enumerate}[leftmargin=2em]
  \item \textbf{制約生成}：source式から\code{Generated}を作る．
  \item \textbf{制約の閉包}：生成した制約を解き，主要な結果を得る．
  \item \textbf{独立した型付け}：実行可能推論器を定義に含めず，型付け規則だけで作る命題
        \code{Typing}で，
        source programが型を持つことを述べる．
\end{enumerate}

健全性・完全性・主要性は，前節で定義した意味で，この三つの層を結ぶ性質である．

\chapter{制約生成：target，hard，pending}

\section{Generatedという三つ組}

Leanでは一つの式から得られる情報を次の構造にまとめる．

\begin{lstlisting}
structure Generated where
  target  : Ty
  hard    : List Equation
  pending : List CheckObligation
\end{lstlisting}

それぞれの意味は次のとおりである．

\begin{definitionbox}{target（生の結果型）}
\code{target}は，式自身を構文に沿って調べた直後の結果型である．
まだ代入を適用していないため，型変数を含んでよい．
「最終的に確定した型」ではなく，制約blockが返そうとしている型の入口である．
\end{definitionbox}

\begin{definitionbox}{hard（無条件に必要な等式）}
\code{hard}は，checking変換の種類に関係なく必ず成り立つ型等式の列である．
たとえば関数適用$f\;x$では，$f$の型が「引数型から結果型への関数」でなければならない．
この外形はhard等式に入る．
\end{definitionbox}

\begin{definitionbox}{pending（後で判定するchecking要求）}
\code{pending}は，ある生の型が期待型として使えるかを後で判定する要求の列である．
本書では$s\checks e$と書き，「source型$s$をexpected型$e$として使えるか」と読む．
普通の型等式になる場合もあれば，matcher producerからslotへの特殊変換になる場合もある．
\end{definitionbox}

\section{関数適用を一つ生成してみる}

関数式$f$から
\[
  \langle t_f,H_f,P_f\rangle
\]
が，引数式$x$から
\[
  \langle t_x,H_x,P_x\rangle
\]
が得られたとする．新しい型変数$d,r$を割り当てると，$f\;x$の生成結果は概略
\[
\begin{aligned}
  \mathit{target}  &= r,\\
  \mathit{hard}    &= H_f \mathbin{+\!+} H_x
                       \mathbin{+\!+} [t_f = d\fn r],\\
  \mathit{pending} &= P_f \mathbin{+\!+} P_x
                       \mathbin{+\!+} [t_x\checks d]
\end{aligned}
\]
となる．ここで$+\!+$は列の連結である．

関数の外形$t_f=d\fn r$は常に必要なのでhardである．一方，引数$t_x$を$d$として
使う方法は，単なる一致かmatcher変換かがまだ分からないのでpendingへ置く．

\begin{pointbox}{「$r$ checks against $d$」の読み方}
$[r\checks d]$は「式の生の型$r$を，使用箇所が要求する型$d$として使えるか後で調べる」
という記録である．$r=d$と最初から決める記号ではない．
同様に$[p\checks d]$は，source型が$p$である別の式に同じ期待型$d$が課されたことを表す．
\end{pointbox}

\section{なぜpendingをすぐ等式にしないのか}

matcher producerとslotの間には，capabilityを調べる特殊変換がある．
したがって「source型をexpected型として使える」という要求は，常に型等式になるとは限らない．
説明のため，次の三つの名前と型がcontextに登録済みだと仮定し，末尾の二つの関数適用を比べる．
最初の三行は説明用の型の前提であり，sourceプログラム本体は最後の二行である．

\begin{lstlisting}
inc : Int -> Int
useMatcher : Slot Any Int -> Bool
m : Matcher Any Int

inc 0
useMatcher m
\end{lstlisting}

ここで\code{Any}は「特定の分解能力を要求しない」capabilityである．\code{m}は実際に作られた
matcherなのでproducer型\code{Matcher Any Int}を持ち，\code{useMatcher}はmatcherを受け取る
場所なのでslot型\code{Slot Any Int}を引数に要求する．

どちらの適用でも，制約生成時には引数位置の型として新しい変数$d$を置き，引数について
「引数の生の型を$d$として使えるか」というpendingを作る．違いが分かるのはhardを解いた後である．

\begin{enumerate}[leftmargin=2em]
  \item \code{inc 0}では，関数側のhard

        \[
          \TyInt\fn\TyInt = d\fn r
        \]

        を解くと$d=\TyInt$になる．したがってpendingは
        $\TyInt\checks\TyInt$となり，普通の等式$\TyInt=\TyInt$として処理できる．
  \pagebreak[3]
  \item \code{useMatcher m}では，関数側のhard

        \[
          (\Slot\,\mathsf{Any}\,\TyInt)\fn\mathsf{Bool} = d\fn r
        \]

        を解くと$d=\Slot\,\mathsf{Any}\,\TyInt$になる．pendingは

        \[
          \Matcher\,\mathsf{Any}\,\TyInt
          \checks
          \Slot\,\mathsf{Any}\,\TyInt
        \]

        となる．これはproducer型とslot型の外側が異なるため型等式ではないが，
        matcherが調べる対象値の型\code{Int}とcapabilityの\code{Any}が適合するので，matcherからslotへの
        特殊変換として受理できる．
\end{enumerate}

もし生成直後にすべてのpendingを等式へ変えると，二番目は
\[
  \Matcher\,\mathsf{Any}\,\TyInt
  =
  \Slot\,\mathsf{Any}\,\TyInt
\]
を要求してしまう．\code{Matcher}と\code{Slot}は型の最外形を作る異なる記号
（型constructor）なので，これは解けず，
本来受理すべき\code{useMatcher m}を誤って拒否する．反対に，すべてをmatcher変換として扱えば，
一番目の普通の\code{Int}引数を処理できない．

生成時点では$d$がまだ新しい型変数であり，expected型の外側が\code{Int}なのか\code{Slot}なのか
確定していない．実際の関数式はさらに別の関数適用や\code{let}を含むこともあり，その場合は
hardをまとめて解くまで外側の形が分からない．そこで現行設計は，まずhardを最も一般的に解き，
その結果をpendingへ適用してから，普通の等式かmatcher変換かを分類する．

pendingは制約を無視する仕組みではなく，\emph{どのchecking規則を使うかだけを延期する仕組み}である．
分類後に得られる型とcapabilityの等式は，最終的にすべて解かなければならない．

この「hardを解く前に型の外側を推測しない」方針を本書では\emph{no-guess}と呼ぶ．

これはあらゆる意味的な解を探索する方式ではない．hardの主要解でchecking分岐を固定した後に
受理できるものを\emph{branch-stable acceptance}と呼ぶ．branch-stableとは，後から型を具体化しても
checkingの分岐選択をやり直さない，という意味である．

\section{小さな例：\code{((lambda z. x) 0)}}

外側の変数$x$の型を$p$とする．式
\[
  ((\lambda z.\,x)\;0)
\]
を生成する．$z$の型を$a$，適用が期待する引数型を$d$，結果型を$r$とすると，
\[
\begin{aligned}
  \mathit{target}  &= r,\\
  \mathit{hard}    &= [a\fn p = d\fn r],\\
  \mathit{pending} &= [\TyInt\checks d]
\end{aligned}
\]
となる．hardを解くと$a=d$かつ$p=r$である．その後pendingを調べると$d=\TyInt$になる．
したがって式全体の結果は$p$，つまり$x$と同じ型である．

ここで$r$は「valueの型」そのものを最初から意味していたのではない．
$r$は右辺式の\code{target}として新しく割り当てた仮の結果変数であり，
制約を解いた後に$p$と同じだと分かる．

\chapter{多相letは何をしているのか}

\section{valueとはletの右辺のこと}

Leanの\code{Expr.letE value body}における\code{value}は，実行時の値一般を指す言葉ではない．
次のsource構文で等号の右側にある式$e_1$を指すフィールド名である．
\[
  \mathbf{let}\ y=e_1\ \mathbf{in}\ e_2
\]
このとき\code{value}は$e_1$，\code{body}は$e_2$である．

\section{full-cut方式}

現行実装の\code{let}はfull-cut方式である．cutは「境界で切る」という意味であり，
右辺の制約blockをbodyより先に完全に閉じる．手順は次のとおりである．

\begin{enumerate}[leftmargin=2em]
  \item 外側contextで右辺\code{value}を生成する．
  \item 右辺のhardを解き，pendingを分類・解決して，ほかの解の元になる最も一般的な
        closure（主要closure）を得る．
  \item closureの代入を外側contextへ適用する．
  \item closureの結果型を，更新後contextに自由でない変数について量化し，型schemeを作る．
        この量化処理を一般化と呼ぶ．
  \item そのschemeをcontext先頭へ追加してbodyを生成する．
  \item 右辺の内部制約は親blockへ出さず，外側contextへ見せる必要がある影響だけを
        有限個の等式（interface equations）として返す．
\end{enumerate}

この実行可能な手順は\code{Source/Elaboration.lean}の\code{elaborate}に現れる．
一方，独立した関係\code{Elaborates.letE}は実行可能推論器を呼ばず，
任意の主要closureとその吸収性の証明を受け取る．

\section{PrincipalBlockClosure}

\code{PrincipalBlockClosure}は，最も一般的な制約blockの解を表す記録である．
ここで「最も一般的」とは，ほかの解をこの解への追加の代入で得られることをいう．
この記録は二段階の解をまとめる．

\begin{enumerate}[leftmargin=2em]
  \item pendingからhardへ安全に移せる要求がなくなるまで処理する．この反復を飽和と呼び，
        飽和後のhardを解く最も一般的な代入を記録する．
  \item まだpendingに残った要求を分類して得る等式（残余等式）を解く，最も一般的な代入を記録する．
\end{enumerate}

二つの代入を順に適用できる一つの代入へまとめたものがclosure substitutionであり，
それを\code{target}へ適用した型がclosure targetである．\emph{吸収的}であるとは，
後から同じ制約を解く代入を合成しても，
選んだ主要代表が正規化された形で保たれる性質である．

「最も一般的」であっても，変数名の代表が文字どおり一意になるわけではない．
等式$p=r$は$r$を$p$へ消しても，$p$を$r$へ消しても解ける．
この小さな自由度が，\code{let}完全性証明の中心問題になる．

\section{一般化とscheme}

一般化は，結果型に現れる変数のうち，更新後contextに自由に現れないものだけを量化する処理である．
その結果である型schemeは，一つの型を複数の具体的な型として再利用できる多相型である．
scheme内の量化変数は名前ではなく位置で表す．これをbound-index scheme，つまり
「量化変数を番号で指すscheme」と呼ぶ．
たとえば恒等関数の型$\alpha\fn\alpha$を空contextで一般化すれば
$\forall\alpha.\alpha\fn\alpha$になる．

一方，右辺の結果型が外側変数$x$の型$p$なら，$p$はcontextにも現れるため量化しない．
\code{let y = ...}で得た$y$はその外側$x$と同じ単相型$p$を持つ．

\section{具体例を最後まで追う}

次の式を考える．

\begin{lstlisting}
lambda x.
  let y = ((lambda z. x) 0) in
  ((lambda w. 0) y)
\end{lstlisting}

前章の計算から，let右辺の結果型は外側$x$の型$p$である．右辺を閉じても$p$は外側contextに
自由に現れるため，$y$のschemeは量化されず，単相$p$になる．bodyの
\code{(lambda w. 0)}は$q\fn\TyInt$という型を持ち，$y:p$へ適用することで$q=p$となる．
bodyの結果型は\code{Int}である．したがってlambda全体の型は
\[
  p\fn\TyInt
\]
である．top-levelで$p$を一般化して読むなら
\[
  \forall p.\ p\fn\TyInt
\]
である．

\begin{warningbox}{\(\forall a.\TyInt\)ではなく\(\forall a.a\fn\TyInt\)}
全体は\code{lambda x. ...}なので，引数$x$を受け取る関数である．bodyが常に0を返しても，
lambdaの矢印は消えない．したがって主要型は$\forall a.a\fn\TyInt$であり，
$\forall a.\TyInt$ではない．後者は実質\code{Int}と同じで，関数でない．
\end{warningbox}

\section{interface equations}

右辺を閉じた代入$S$が外側contextの自由変数を変更した場合，その影響をbodyより外へ伝える
必要がある．ただし右辺の全hardとpendingを再び外へ出すとfull-cutではなくなる．
そこで，contextの各自由型変数$x$について
\[
  x = S(x)
\]
だけを作る．capability変数についても同様である．これが
\code{Context.interfaceEquations}である．

親の後続代入$T$がこのinterfaceを解くことは，contextに見える各変数について
\[
  T(x)=T(S(x))
\]
が成り立つことと同値である．interfaceは右辺の制約全部でも，bodyへ渡すchecking要求でもない．
外側から観測できる有限な効果だけである．

最後に\code{Generated.fromLet effects body}は
\[
\begin{aligned}
  \mathit{target} &= \mathit{body.target},\\
  \mathit{hard} &= \mathit{effects}\mathbin{+\!+}\mathit{body.hard},\\
  \mathit{pending} &= \mathit{body.pending}
\end{aligned}
\]
を返す．右辺valueのpendingは境界を越えないが，body自身が作ったpendingは外側blockに残る．

\chapter{なぜletの完全性証明が難しかったのか}

\section{実行版は一つを選び，関係版は任意の代表を許す}

公開推論器は実行可能な\code{unify}を使い，右辺に対して一つのclosureを計算する．
しかし関係\code{Elaborates.letE}は，条件を満たす任意の主要かつ吸収的なclosureを許す．
完全性を示すには，関係側がどの正しい代表を選んでも，公開推論器が選ぶ代表へ運べることを
示さなければならない．

先ほどの右辺ではhardから$p=r$が得られた．次の二つはどちらも正しい主要な向きである．
\[
  S_L=[r\mapsto p],\qquad S_R=[p\mapsto r].
\]
$S_L$では閉じたcontextは$x:p$，右辺結果も$p$である．
$S_R$では閉じたcontextは$x:r$，右辺結果も$r$である．
意味は同じだが，bodyの生成に使う変数名が違う．interfaceも一方では$p=p$という自明式，
他方では$p=r$という非自明な別名等式になり得る．別名等式とは，異なる変数名が
同じ型を表すことを記録する等式であり，以下ではalias equationとも呼ぶ．

\section{単純な大域的renameでは直らない}

最初に考えたくなるのは「一方の変数名を全体で置き換えればよい」という案である．
本書では，型変数とcapability変数の名前を一対一に付け替える操作をrenameと呼ぶ．
しかし全単射な変数名変更は，自明式$x=x$を別の自明式$y=y$へ移すことしかできない．
非自明式$x=y$へは変えられない．現行リポジトリにはこの失敗を
\code{GlobalRenamingCounterexample}としてLeanで保存している．

ほかにも次の強すぎる案が反例で退けられた．

\begin{itemize}[leftmargin=2em]
  \item 常に左側だけへ別名等式を足す．実際には代表選択により足す側が変わる．
  \item hardだけを共通化し，pendingは文字どおり同じだと要求する．bodyのpendingにも代表名の差が残る．
  \item bodyだけを先にrenameし，任意の外側frameで同値だとする．外側制約が変数の同一性を観測できる．
\end{itemize}

ここでframeとは，生成済み制約blockを穴へ入れる周囲の制約構造である．
完全性の帰納法で再利用するには，裸のblockだけでなく適切な周囲へ入れても受理を保つ必要がある．

\section{必要なのは文字列一致ではなく意味の一致}

二つのpendingが別の変数名を含んでも，周囲のhardを解くすべての代入の下で同じ要求になればよい．
これが\code{EntailedObligationEq H left right}である．\emph{entailed}は
「hard理論$H$から必ず従う」という意味である．

たとえばhardに$p=r$があれば
\[
  p\checks d
  \quad\text{と}\quad
  r\checks d
\]
は文字どおり違うが，$H$を解く任意の代入の下では同じになる．
この関係は，checking分岐の分類，hardへの昇格，飽和，残余等式，最終受理まで輸送できる．

\chapter{M2の最終解法：四つの層}

最終証明は，一つの強いrename定理で全部を済ませない．役割の違う四層を組み合わせる．
以下でsupportとは，型・制約・代入などに実際に現れる有限個の変数名の集合をいう．
freshnessとは，新しく割り当てた変数名が既存のsupportと衝突しない性質をいう．

\begin{center}
\begin{tikzpicture}[
  node distance=7mm and 7mm,
  box/.style={draw,rounded corners,align=center,minimum width=36mm,minimum height=11mm,fill=blue!3},
  arr/.style={-{Latex[length=2mm]},thick}
]
\node[box] (value) {二つのvalue導出\\二つの主要closure};
\node[box,right=of value] (rename) {body renaming\\同じcontextへ輸送};
\node[box,below=of value] (graph) {visible alias graph\\interfaceを共通化};
\node[box,right=of graph] (entailed) {entailed alignment\\target/pendingを意味比較};
\node[box,below=12mm of $(graph)!0.5!(entailed)$,minimum width=55mm] (whole) {whole let certificate\\受理と主要結果を保存};
\draw[arr] (value) -- (rename);
\draw[arr] (value) -- (graph);
\draw[arr] (rename) -- (entailed);
\draw[arr] (graph) -- (entailed);
\draw[arr] (graph) -- (whole);
\draw[arr] (entailed) -- (whole);
\end{tikzpicture}
\end{center}

\section{第1層：entailed alignment}

\code{EntailedGeneratedAlignment}は，二つの\code{Generated}について次を要求する．

\begin{enumerate}[leftmargin=2em]
  \item hardの解集合が同じである．
  \item そのhardを解く任意の代入の下でtargetが同じになる．
  \item 対応するpendingが同じ要求になる．
\end{enumerate}

さらに\code{SupportedEntailedAlignmentCertificate}は，終了supply，左右に追加する有限な
別名等式，それらが既存名と衝突しないことを示すfreshness，frameのsupportをまとめる．
lambda，application，tuple，constructor，primitive，conditionalは，この証明書を構文に沿って
合成できる．

\section{第2層：bodyを同じcontextへ運ぶ有限rename}

二つのclosureの結果contextは$x:p$と$x:r$のように文字どおり異なり得る．
bodyの帰納仮定を使うには，contextを同じにしなければならない．
そこでclosureの有限supportから変数対応$\rho$を作る．

$\rho$には二つの重要な性質がある．

\begin{description}[leftmargin=9em,style=nextline]
  \item[bodyContext\_exact]
    左のclosed contextと一般化schemeを$\rho$で移すと，右のものと文字どおり一致する．
  \item[FixesAtOrAbove]
    value終了後に新しく割り当てる将来の変数名を$\rho$が変更しない．
\end{description}

後者により，左bodyの導出を同じsource AST，同じ数値supplyのまま右contextへ輸送できる．
これはsourceのde Bruijn indexを変える操作ではない．型変数とcapability変数の有限な名前変更である．

\section{第3層：visible closure graph}

bodyを運ぶrenameだけでは，外側\code{fromLet}のinterface差を消せない．
そこで二closureの変数対応のうち，closed contextから実際に観測でき，かつ名前が動く辺だけを
有限個の「左の変数名と右の変数名の対応」として取り出す．この有限な対応表を
visible closure graph，つまり外側contextから見えるclosure変数対応と呼ぶ．

必要な側へfreshな別名等式を追加すると，左右のinterfaceは共通のhard理論
\[
  I_L\mathbin{+\!+}I_R
\]
を表すようになる．別名等式は左右非対称でよい．一つの全体swapを要求しない．

先ほどの例で，一方が$p=p$，他方が$p=r$なら，欠けている側へだけfreshな別名を追加して
共通理論へ移す．ただしcontextやbodyで既に使われている端点をfreshだと呼んではならない．
そのためgraphは全closure supportではなく，contextから見える移動辺に限定する．

\section{第4層：well-formed supplyとfreshness}

freshnessは「この変数名はたまたま違う」という主張ではない．source生成変数の由来と
数値区間を使って証明する．開始supplyがcontextに対してwell-formedなら，sourceが新しく作る変数は
半開区間$[\mathit{start},\mathit{finish})$にあり，古いcontext変数と衝突しない．

closureの局所性により，block supportにない変数はclosure substitutionが変更しない．
これらを合わせると次が言える．

\begin{itemize}[leftmargin=2em]
  \item graph aliasのfresh端点はvalueの割当区間内にある．
  \item 対応する側のinterfaceとbody contextにはその端点が現れない．
  \item value区間のaliasと，後続bodyが作るaliasは衝突しない．
  \item body用renameはbody以後のfresh名を固定する．
\end{itemize}

\section{whole letの組み立て}

bodyの証明書はrename後の座標で作られている．そのため左body側のalias列を$\rho$の逆向きに
たどり，元の座標へpull backする．pull backとは，移した名前を元の名前空間へ戻す操作である．

最後に，左右のgraph alias，pull backしたbody alias，bodyのentailed alignmentをまとめて，
元の二つの\code{Generated.fromLet}全体に対する
\code{SupportedEntailedAlignmentCertificate}を構成する．

この手順が重要なのは，反例で偽になった「孤立したbody renameは任意の周囲で無害」という
中間命題を仮定しない点である．value導出，closure，transport済みbody導出を保持したまま，
interfaceとbodyを一緒に組み立てる．

\chapter{定理が公開推論へつながるまで}

\section{証明の主な鎖}

現行の無条件なM2証明入口は\code{Source/FullM2Completion.lean}に集約されている．
大きな流れは次のとおりである．

\begin{enumerate}[leftmargin=2em]
  \item \code{fullM2LetGraphAliasPresentationComplete}がvisible graph presentationを構成する．
  \item \code{fullM2LetSupportedAssemblyBridge}がwhole-\code{let}証明書を構成する．
  \item \code{fullM2CoherenceComplete}が通常構文の帰納法と\code{let}の場合を統合する．
  \item \code{FullM2Replay}が任意の関係的導出を実行可能な導出へ再生する．
  \item \code{wellFormedElaborationPrincipalityComplete}が完全性と主要性をまとめる．
\end{enumerate}

ここで\emph{coherence}は，同じsourceに対する異なる正しい導出が，受理と主要結果について
整合することをいう．\emph{replay}は，関係的導出を入力として，公開推論器が実際に選ぶclosureを
使った導出を再構成することである．

\section{公開API}

最終的に次の定理が公開される．

\begin{longtable}{p{58mm}p{82mm}}
\toprule
定理 & 初学者向けの読み方 \\
\midrule
\endfirsthead
\toprule
定理 & 初学者向けの読み方 \\
\midrule
\endhead
\path{Typing.infer_isSome} & 独立した型付けがあれば，公開推論は失敗しない．\\
\path{Inference.typable_iff_infer_isSome} & well-formed signatureの下で，型が付くことと推論成功が同値．\\
\path{Inference.typableDecidable} & 型が付くかどうかを公開推論で決定できる．\\
\path{Inference.infer_success_principalResult} & 推論が返した型は，すべての型付け結果の主要型である．\\
\path{Inference.infer_principalResult} & 関係側の主要な型付けから，公開推論の主要結果を得られる．\\
\path{PrincipalTyping.finiteRenamingEq} & 二つの主要代表は，有限個の変数名変更を除いて一致する．\\
\bottomrule
\end{longtable}

\code{Typing.infer\_isSome}と\code{finiteRenamingEq}はsignatureのwell-formednessを仮定しない．
推論成功から関係的型付けを得る健全性まで使うAPIは，\code{Signature.WellFormed}を明示的に受け取る．
signatureがwell-formedであるとは，名前が重複せず，各型schemeが外部の未定義変数を含まず，
constructorの引数数や結果のdata型が宣言表と一致することである．

\section{負例にも完全性が効く}

\code{M2Regression.cutRejectsBackflow}は，\code{let}右辺を閉じた後に，bodyから互換でない要求を
右辺へ逆流させようとする例である．実行可能推論が\code{none}を返すだけなら，推論器が不完全で
見逃した可能性が残る．M2完全性により，もし独立した\code{Typing}が存在すれば推論は成功するはずである．
実際には失敗するので，\code{cutRejectsBackflow\_not\_typable}として非型付けまで証明できる．

\chapter{M3からM5までの現在地}

\section{M3：宣言と通常の呼出し}

M3ではBool/Listのdata型，data constructor，pattern constructor，primitive，有限signatureを追加した．
data constructorは\code{nil}や\code{cons}のようにdata値を作る名前であり，primitiveは
\code{add}や\code{map}のように実装があらかじめ与えられた組込み演算である．
constructorとprimitiveの引数処理は，通常の関数適用\code{Generated.fromApp}を左から畳み込む．
したがってM2で証明した通常applicationの構造を再利用できる．

\code{add}，\code{append}，\code{member}，\code{deleteFirst}，\code{map}について，
実装が正しいものとして固定する標準の型scheme（canonical scheme），
\code{ifE}，名前・引数数・型不一致の正負回帰は実装済みである．
M3が表で「一部」なのは，M4のpattern/matcherと合わせた論文listing全体の静的回帰が残るためである．

\section{M4：patternとmatcher}

source構文は\code{fixE}，\code{matcher}，\code{matchAll}，\code{matchFirst}，
value pattern，pattern constructor，conjunction pattern，pattern function applicationまで持つ．
conjunction patternとは，同じ対象へ二つのpatternを左から右の順に適用するpatternである．
M4の実行可能elaborationと関係的elaborationは同じ構文を再帰的に処理し，
成功から\code{M4.Typing}を得る健全性を証明済みである．
ここでelaborationはsource式から型制約を生成する処理である．実行可能版はLeanの関数として
結果を計算し，関係版は「この生成結果が規則から導ける」という命題として同じ処理を表す．

fresh変数の由来を追うsupport provenanceも証明している．supportとは，型や制約に
実際に現れる変数の有限集合である．provenanceは「由来」という意味であり，ここでは各変数が
もとのcontextにあったか，開始supplyから終了supplyまでの間に新しく割り当てられたことを表す．
well-formed signature，つまり同じ名前の宣言が重複せず，各schemeが正しく閉じている宣言表の下では，
最終的な相互再帰関係\code{ElaboratesFuel}と，公開関係\code{M4.Elaborates}を含む全M4構文について
この性質を証明済みである．公開されたsupply増加と由来証明をまとめた\code{M4.supplyAndSupport}も完成した．
これは生成変数の管理が正しいという結果であり，M4全域の完全性が得られたという意味ではない．

一方，M2--M3で完了した「型が付くなら必ず推論成功」と「返却型の主要性」はM4全域では未証明である．
matcher literalを構文だけから検査する処理（static check），\code{fix}直下matcherの型形状，clause内部のslot要求，
および\code{let}での任意closure選択を同時に扱う必要があるため，独立した大きな課題である．

\section{M5：評価と型安全性}

M5にはプログラム実行中に実際に得られる値（runtime value），関数本体と定義時の環境を組にした
関数closure，matcher本体と環境を組にしたmatcher closure，候補を記録順に奥まで調べる深さ優先探索，
\code{matchAll}と\code{matchFirst}の関係的評価とfuel付き評価器がある．
fuelは再帰計算を進められる回数の上限である．有限な成功導出について，実行可能評価の健全性，
完全性，fuel単調性を証明している．fuel単調性とは，あるfuel量で成功した計算が，より大きい
fuel量でも同じ結果で成功することである．

conjunction patternは組込み規則A-ANDで実行する．A-ANDは一つの照合課題を，同じmatcherと同じ対象を持つ
左pattern，右patternの二課題へこの順で置き換える．左側で得たbindingが右側のvalue patternから見えるため，
静的型付けと実行の変数順が一致する．

\section{一般pattern functionとMNode}

pattern functionはpatternを引数に取り，別のpatternを組み立てる機能である．本体内で一時的に束縛した
変数は外側の\code{matchAll}へ見せてはならない一方，引数として渡されたpatternが作る束縛は外側へ返す必要がある．
この二つを区別するため，実行時にはMNodeを使う．MNodeは，pattern function本体専用の照合課題，環境，
private binding（本体の中だけで使う束縛）を持つ隔離された探索nodeである．

具体例はPaper 2の\code{pair}である．

\begin{lstlisting}
def pattern pair (first, rest) :=
  ($private & ~first) :: #private :: ~rest
\end{lstlisting}

\code{\$private}は最初の要素を本体内だけで束縛し，\code{\#private}が次の要素と同じかを検査する．
\code{\textasciitilde first}と\code{\textasciitilde rest}は，呼出し側から渡された第一・第二引数patternを表す．MNodeの処理は次の順である．

\begin{enumerate}[leftmargin=2em]
  \item \code{pair q1 q2}をmatcherへ渡す前に，\code{pair}本体を持つMNodeへ展開する．
  \item 本体の通常の照合はMNode内部で進め，\code{private}の束縛も内部に保存する．
  \item \code{\textasciitilde first}または\code{\textasciitilde rest}へ着いたら，対応する実引数patternを外側の探索へ戻す．
  \item 本体が終わったらMNode内部の束縛を捨て，実引数patternが作った外側の束縛だけを残す．
\end{enumerate}

また，pattern function適用はuser-defined matcherのcatch-all clauseより先に処理する．catch-all clauseとは，
それ以前のclauseに一致しなかったpatternを一般的に受け取る最後のclauseである．先にMNodeへ展開しなければ，
\code{pair}そのものがcatch-allに捕まり，本体が実行されないからである．

正確な7-clause source \code{multiset} matcherを使う回帰では，

\begin{lstlisting}
matchAll [1,2,1,3] as multiset integer with
  $outer :: pair $inner _ -> [outer, inner]
\end{lstlisting}

の照合結果に対応するbinding列として，

\[
 [[1,2],[1,2],[1,3],[1,3]]
\]

を外側，内側の順で並べ，順序と重複を保って得る全\code{matchAll}式をLean kernelで確認している．
同じ探索を使う\code{matchFirst}が先頭の\code{[1,2]}を返すことも固定した．\code{private}は各結果に含まれない．
pattern function本体の静的な検査証明も，保存schemeの標準的な一つの具体化について，生成された等式とchecking条件を実際に満たす代入と，代入後の結果型の一致を保持する．これは，全具体化に対するinterface一致や主要性の証明ではない．
さらに，source/runtime定義表の双方向対応を\code{PatternFunctionDefinitions.Agree}で表す．これは，
runtime表の各本体が上記の標準的な具体化に対するsource側の検査証拠を持ち，source側の各宣言にruntime実装があるという条件である．
したがって，runtimeだけにある未検査本体と，実装のないsource宣言の両方を除ける．\code{pair}の定義表が
この条件を満たすことも証明済みである．

MNode探索は\code{evalCheckedPatternFunctionNodesFuel}という公開fuel評価器にも接続した．この評価器は定義表と
上記の双方向対応の証明を受け取り，全\code{Expr}を再帰的に評価し，\code{matchAll}と\code{matchFirst}ではMNode探索を使う．成功した
\code{matchAll}の探索部分について，実行可能な一歩関数を使う帰納的な順序付き深さ優先導出も証明済みである．ただし評価器全体に
対応する独立な関係的評価\code{Eval}と，MNodeを含む型安全性はまだ証明していない．

型安全性は整数，真偽値，List，tuple，型付き環境，通常・再帰closure，任意の型付き関数位置を持つapplication，
\code{map}，および明示的なcallback仮定の下のbuilt-in matcher断片までである．
callbackとは，処理本体を引数として受け取り，必要な時点で呼び出す関数である．built-in matcherとは，
利用者がsourceで定義するのではなく，実行器に直接実装されたmatcherである．
ここで「任意の型付き関数位置」とは，applicationの左側が構文上直接のlambdaや\code{fixE}でなくても，
実行時型付けで関数型を持てばよいという意味である．例えば条件分岐が返したclosureの適用も扱える．
sourceの多相型付けからこの関数断片を再構成する橋，user-defined matcher clause，
一般pattern functionのMNodeを同じ型安全性契約へ接続する作業が残る．ここで保存とは，実行の一歩後も型が保たれる
ことであり，進行とは，型の付いた実行可能状態が結果でなければ次の一歩を進められることである．
残作業はこの保存・進行の形を持つが，
多相\code{let}のschemeとruntime環境を結ぶ設計を含むため，小さな機械作業ではない．

\chapter{multisetWithListArgumentを理解する}

\section{Paper 1のmultiset定義との対応}

現行source AST，つまりsource式の構文を木として表すデータ型には，プログラム最外部
（top-level）の利用者定義値を名前で参照するmodule環境がない．module環境とは，
名前からその定義済みの値を引く表である．
通常の値参照はde Bruijn indexで表す．de Bruijn indexとは，変数名の代わりに
「最も内側のbinderから何個外か」を自然数で表す方法である．

Paper 1の概略
\begin{lstlisting}
def multiset m := matcher { ... uses list m ... }
\end{lstlisting}
では，\code{list}は既に定義済みのlibrary matcherである．Lean fixtureでは
\code{multisetDefinition}の外側環境に\code{list}が一つあるものとして，内部からindex 2で参照する．
index 0は\code{fixE}の引数$m$，index 1は再帰的な自己\code{multiset}，index 2が外側の\code{list}である．

\section{三つの式の役割}

\begin{lstlisting}
multisetDefinition
  = fixE (matcher multisetClauses)

multisetWithListArgument
  = lambda list. multisetDefinition

closedMultisetDefinition
  = multisetWithListArgument listMatcherDefinition
\end{lstlisting}

\code{multisetDefinition}自体が，Paper 1の再帰的
\code{def multiset m := ...}に対応する．\code{fixE}を適用すると，変数から実行中の値を引く表
（実行時環境）は
\code{argument :: self :: definitionEnvironment}となる．

\code{multisetWithListArgument}はPaper 1に表示される別のアルゴリズムではない．
自由依存である\code{list}をlambda引数へ持ち上げた，closure conversion風の補助式である．
closure conversionとは，関数が外側から参照する自由な値を，明示的な環境や引数として渡す変換である．

\code{closedMultisetDefinition}はsourceで定義した\code{listMatcherDefinition}を実際に渡し，
空のsource contextから使える閉じた式にする．閉じた式とは，外側の変数を参照せず，それだけで
評価できる式である．静的確認だけなら，\code{multisetDefinition}を
既知のlist schemeを持つcontextで調べてもよい．end-to-endの実行回帰では隠れたprimitive値を使わず，
空context・空の実行時環境（runtime environment）から始めるため，この明示的な適用を置いている．

\section{なぜlambda版単独は現在拒否されるのか}

\code{multisetWithListArgument}の制約生成自体は成功し，hard制約も解ける．
失敗は残ったchecking要求の解決で起きる．外側lambda引数\code{list}の型は概略
\[
  \Slot\,\chi\,\alpha \fn \beta
\]
までhardから分かるが，返り型$\beta$がmatcher producerであることはhardだけから決まらない．

clause内には概略
\[
  \Prod[\beta,\ \Matcher(\mathsf{List}\,\chi)(\mathsf{List}\,\alpha)]
  \checks \text{product-slot}
\]
という要求が残る．productからproduct-slotへの特殊変換を選ぶには，productの全成分が
構文上すでにmatcherだと分かる必要がある．$\beta$が変数のままなので，resolverは普通の等式へ進み，
最外constructorの不一致で失敗する．

\begin{pointbox}{型が存在しないのではない}
$\beta$を適切なmatcher型へ具体化すれば，すべてのchecking要求を満たす意味的な解はある．
期待される全体型は概略
\[
  (\Slot\,\chi\,\alpha\fn
    \Matcher(\mathsf{List}\,\chi)(\mathsf{List}\,\alpha))
  \fn
  \Slot\,\chi\,\alpha\fn
    \Matcher(\mathsf{List}\,\chi)(\mathsf{List}\,\alpha)
\]
である．しかしその形を選ぶには，hardの主要解に現れていないmatcher外形を$\beta$へ先回りして与える
必要がある．現在のbranch-stable/no-guess仕様はこの推測をしない．
\end{pointbox}

\code{closedMultisetDefinition}では，実引数\code{listMatcherDefinition}の既知の型が外側lambdaの
domainへhardな情報を与え，$\beta$がmatcherだと確定する．その後なら同じproduct要求を特殊変換として
分類でき，推論が成功する．

この現象は偶然のM4実装バグではなく，\code{GeneratedSemanticAcceptanceRegression}にある
M1の最小境界の大規模な実例である．「完全に具体化すれば意味的にchecking可能」と
「hardの主要解で分岐を固定した公開手続きが受理する」は同じではない．

\section{設計上の選択肢}

この高階matcher constructorをsource言語でどこまで受理したいかにより，選択肢が変わる．

\begin{enumerate}[leftmargin=2em]
  \item \textbf{既知型contextまたは明示注釈を使う}．no-guessを保つ最も安全な案である．
  \item \textbf{matcher位置から期待slotを逆向きに伝える}．双方向型付けに近い規則を追加する．
        どの構文位置が十分な根拠になるかを明示する必要がある．
  \item \textbf{checking分岐を探索する}．より多くの意味的解を受理できるが，
        branch-stable設計と主要性の証明を大きく変更する．
\end{enumerate}

現状では，Paper 1の\code{multiset}本体そのものが無意味なのではない．
名前付きtop-level環境を持たないfixtureを閉じるための外側lambdaが，
高階matcher constructorに対する既知のno-guess境界を露出している，という位置づけである．

\section{別問題：kernelで大きな推論を計算すること}

大きいmatcherの検証例（fixture）では\code{M4.infer}の実行結果を\code{\#eval}で確認できても，
Lean kernelの定義簡約だけで同じ等式を\code{rfl}により証明することが難しい．Lean kernelとは，
作られた証明がLeanの規則に従うかを最後に確認する，小さな検査器である．
測度に基づく再帰で定義したpattern検査や単一化が簡約を遮るためである．

これは前節のbranch-stable拒否とは別問題である．改善するなら，既存のfuel構造再帰
\code{M4.elaborateFuel}をsolverまで引数化し，pattern検査・direct-self検査・単一化にも
計算可能なfuel版を用意する．その成功を公開版へ輸送する定理を証明した後，
小さいcatch-all，list matcher，closed multisetの順でkernel回帰を固定するのが自然である．

\chapter{コードを読むための案内}

\section{最初に読むファイル}

\begin{longtable}{p{62mm}p{78mm}}
\toprule
ファイル & 読める内容 \\
\midrule
\endfirsthead
\toprule
ファイル & 読める内容 \\
\midrule
\endhead
\path{TypePM/Constraints.lean} & \code{Equation}，\code{CheckObligation}，\code{Generated}の最小定義．\\
\path{TypePM/Source/Syntax.lean} & M4までの最終source AST．\\
\path{TypePM/Source/Elaboration.lean} & 実行可能生成と独立した\code{Elaborates}，\code{fromApp}，\code{fromLet}．\\
\path{TypePM/BlockClosure.lean} & 右辺blockを閉じる\code{PrincipalBlockClosure}．\\
\path{TypePM/ContextInterface.lean} & 外側contextへの有限なinterface equations．\\
\path{TypePM/Source/M2Regression.lean} & 多相identity，明示的let，full-cut負例．\\
\path{TypePM/Source/GlobalRenamingCounterexample.lean} & 一つの大域renameでは足りない具体的反例．\\
\path{TypePM/EntailedPendingTransport.lean} & pendingをhard理論の全解の下で比較し，受理を輸送する層．\\
\path{TypePM/Source/EntailedAlignment.lean} & source blockの意味的な比較証明書．\\
\path{TypePM/Source/WholeLetEntailedAssembly.lean} & interfaceとpull backしたbody aliasのwhole-let合成．\\
\path{TypePM/Source/FullM2GraphAliasPresentation.lean} & visible closure graphのsource由来構成．\\
\path{TypePM/Source/FullM2Coherence.lean} & 通常構文とletの場合を統合する帰納法．\\
\path{TypePM/Source/FullM2Replay.lean} & 関係的導出から実行可能導出への再生．\\
\path{TypePM/Source/FullM2Completion.lean} & 完了したM2の公開入口．\\
\path{TypePM/Source/Paper1Programs.lean} & list/multiset matcherの完全なsource fixture．\\
\path{TypePM/GeneratedSemanticAcceptance}\newline
\path{Regression.lean} & branch-stable受理と意味的解の差を示す最小例．\\
\bottomrule
\end{longtable}

\section{検証コマンド}

Lean全体はリポジトリ直下で次のように検証する．

\begin{lstlisting}
lake build
git diff --check
\end{lstlisting}

追加公理への依存は\path{TypePM/AxiomAudit.lean}で検査する．
M2の公開入口はLean標準の\code{propext}，\code{Classical.choice}，\code{Quot.sound}以外の
追加公理に依存しない．

本書は\code{tex}ディレクトリで次のように生成する．

\begin{lstlisting}
make
\end{lstlisting}

出力は\code{type-pm-mech2-guide.pdf}である．

\chapter{まとめと再開地点}

\section{M2について覚えておく四文}

\begin{enumerate}[leftmargin=2em]
  \item \code{let}右辺はfull-cutで先に閉じ，右辺pendingをbodyへ逆流させない．
  \item 任意の主要closureは意味的には同じでも，代表変数とinterfaceが文字どおり違い得る．
  \item bodyはfuture名を固定する有限renameで同じcontextへ運び，outer interfaceは別のvisible alias graphで共通化する．
  \item pendingは文字列一致でなく，hardを解くすべての代入の下で同じかどうかを比較する．
\end{enumerate}

この四層をwell-formed supplyによるfreshnessが支え，公開推論の完全性・受理同値・主要性へつながる．

\section{次の作業はM2ではなくM4/M5}

M2のill-formed開始supplyは仕様外であり，完了のために追加証明する項目ではない．
次の静的な主作業は，大きいPaper 1の検証例をLean kernelで計算可能な形にし，正例の正確な
\code{M4.infer}結果と\code{M4.Typing}，負例の非型付けを固定することである．

fresh変数の由来証明は公開\code{M4.Elaborates}まで完成したため，次はM4全域の完全性・主要性を
別計画として調べる．M5ではsourceの多相\code{let}から，すでに
任意の型付き関数位置と\code{map}を扱える実行時型付けへの橋を作り，user-defined matcher clauseと
pattern functionのMNodeを一般の実行時型付けと，型の不整合による停止が決して起きないという性質
（no-stuck）へ接続する．

\appendix

\chapter{現行実装の型規則一覧}

\section{この一覧の読み方}

現行実装の型付けは，式から最終型を一度に決めるのではなく，次の二段階に分かれる．

\begin{enumerate}[leftmargin=2em]
  \item \textbf{制約生成}：式を左から右へ調べ，\code{target}，\code{hard}，\code{pending}を作る．
  \item \textbf{制約を閉じる処理}：hardを単一化し，pendingのchecking規則を分類して解き，
        最終的な主要型を得る．ここで主要型とは，ほかの解を代入によって得られる最も一般的な型である．
\end{enumerate}

したがって，本付録で中心になるのは通常の「$e$の型は$\tau$である」という判断より一段手前の
\emph{制約生成判断}である．次の記号を使う．

\[
  \Sigma;\Gamma;s \vdash e \Longrightarrow
  \langle t;H;P\rangle;s'
\]

これは「宣言表$\Sigma$とcontext $\Gamma$の下で，開始supply $s$から式$e$を調べると，
target $t$，hard列$H$，pending列$P$を生成し，次の未使用番号が$s'$になる」という意味である．
宣言表とはdata constructor，primitive，pattern constructorなどの型を登録した表である．
$\langle t;H;P\rangle$を以下では生成blockと呼ぶ．$H_1\mathbin{+\!+}H_2$は列を順に連結する記号である．

型変数とcapability変数には別々の番号がある．$s$から新しく取る型変数を$\alpha_s$，
capability変数を$\kappa_s$と略記する．実装上の$s$はこの二種類の次番号を持つ組である．
「freshな変数」とは，それまで使われていない新しい番号の変数を指す．

\begin{pointbox}{規則の上と下}
横線の上は，その規則を使うために先に確かめる条件である．横線の下は，条件がそろったときに
導ける結論である．条件が複数ある場合は，source順に左から右へ調べる．本付録では読みやすさのため，
supplyの単純な加算や列要素の再帰的な連結は一部を文章で表す．省略された検査を意味しない．
\end{pointbox}

\section{生成blockを組み立てる共通操作}

以下の四つを押さえると，個々の規則を短く読める．

\begin{longtable}{p{35mm}p{103mm}}
\toprule
操作 & 生成するblock \\
\midrule
\endfirsthead
\toprule
操作 & 生成するblock \\
\midrule
\endhead
lambda & bodyが$B=\langle t_B;H_B;P_B\rangle$なら，
$\mathsf{Lam}(d,B)=\langle d\fn t_B;H_B;P_B\rangle$．引数型$d$を結果型の前へ付ける．\\
application & 関数$F=\langle t_F;H_F;P_F\rangle$と引数$A=\langle t_A;H_A;P_A\rangle$に対して
\[
\begin{aligned}
\mathsf{App}(F,A;d,r)=\langle r;{}&
H_F\mathbin{+\!+}H_A\mathbin{+\!+}[t_F=d\fn r];\\
&P_F\mathbin{+\!+}P_A\mathbin{+\!+}[t_A\checks d]\rangle .
\end{aligned}
\]
$d$は引数位置の型，$r$は適用結果の型で，どちらもfreshに取る．\\
checked expression & $G=\langle t;H;P\rangle$をexpected型$u$として使う場合，
$\mathsf{Checked}(G,u)=\langle H;P\mathbin{+\!+}[t\checks u]\rangle$．
結果型を返さず，checking要求を一件追加する補助形である．\\
let & bodyが$B=\langle t_B;H_B;P_B\rangle$，右辺を閉じた効果が等式列$I$なら，
$\mathsf{Let}(I,B)=\langle t_B;I\mathbin{+\!+}H_B;P_B\rangle$．
右辺自身のpendingは閉じる処理で解決済みなので外へ出さない．\\
\bottomrule
\end{longtable}

\section{M2--M3の式の規則}

ここでschemeは量化変数を持てる多相型であり，instantiateは量化変数をfreshな自由変数へ
置き換える操作である．以下の規則は\path{TypePM/Source/Elaboration.lean}の
\code{Elaborates}，\code{ElaboratesItems}，\code{ElaboratesCall}に対応する．

\subsection{変数，整数，\code{something}}

\[
\frac{
  \Gamma[i]=\sigma
  \qquad \mathsf{instantiate}(\sigma,s)=(\tau,s')
}{
  \Sigma;\Gamma;s\vdash \mathsf{var}\ i
  \Longrightarrow\langle\tau;[];[]\rangle;s'
}\;\textsc{Var}
\]

contextの$i$番目にあるscheme $\sigma$をその使用場所でinstantiateする．したがって多相変数は
使用するたびに別のfresh変数を得る．

\[
\frac{}{\Sigma;\Gamma;s\vdash n
  \Longrightarrow\langle\TyInt;[];[]\rangle;s}\;\textsc{Int}
\]
\[
\frac{}{\Sigma;\Gamma;s\vdash \mathsf{something}
  \Longrightarrow\langle\Matcher\,\mathsf{Any}\,\alpha_s;[];[]\rangle;s'}
  \;\textsc{Something}
\]

整数はhardもpendingも作らない．\code{something}は任意の対象型$\alpha_s$に対する
matcher producerを作り，型変数を一つ消費する．

\subsection{lambda，application，tuple}

\[
\frac{
  \Sigma;(\alpha_s::\Gamma);s_1\vdash e
  \Longrightarrow B;s'
}{
  \Sigma;\Gamma;s\vdash \lambda e
  \Longrightarrow\mathsf{Lam}(\alpha_s,B);s'
}\;\textsc{Lam}
\]

$s_1$は$\alpha_s$を一つ予約した後のsupplyである．$\alpha_s::\Gamma$は，lambdaの引数型を
contextの先頭へ追加することを表す．

\[
\frac{
  \Sigma;\Gamma;s\vdash f\Longrightarrow F;s_1
  \qquad
  \Sigma;\Gamma;s_1\vdash a\Longrightarrow A;s_2
}{
  \Sigma;\Gamma;s\vdash f\ a
  \Longrightarrow\mathsf{App}(F,A;d,r);s_3
}\;\textsc{App}
\]

$d,r$は$s_2$から取るfreshな型変数である．関数の外形$t_F=d\fn r$はhardへ入るが，
引数の使用可能性$t_A\checks d$はpendingへ入る．第2.3節の具体例はこの規則の二通りの解決である．

\[
\frac{
  \Sigma;\Gamma;s\vdash[e_1,\ldots,e_n]
  \Longrightarrow[(t_1,H_1,P_1),\ldots,(t_n,H_n,P_n)];s'
}{
  \Sigma;\Gamma;s\vdash(e_1,\ldots,e_n)
  \Longrightarrow
  \langle\Prod[t_1,\ldots,t_n];H_1\mathbin{+\!+}\cdots\mathbin{+\!+}H_n;
  P_1\mathbin{+\!+}\cdots\mathbin{+\!+}P_n\rangle;s'
}\;\textsc{Tuple}
\]

要素列は左から右へ調べる．空tupleでは型が$\Prod[]$，hardとpendingが空になる．

\subsection{多相\code{let}}

\[
\frac{
  \begin{gathered}
  \Sigma;\Gamma;s\vdash v\Longrightarrow V;s_v,\qquad
  C\text{は}V\text{の吸収的な主要closure},\\
  \Sigma;(\mathsf{gen}(C(t_V),C(\Gamma))::C(\Gamma));s_b
    \vdash e\Longrightarrow B;s'
  \end{gathered}
}{
  \Sigma;\Gamma;s\vdash \mathsf{let}\ v\ \mathsf{in}\ e
  \Longrightarrow\mathsf{Let}(\mathsf{interface}(\Gamma,C),B);s'
}\;\textsc{Let}
\]

$v$がletの右辺，$e$がbodyである．closure $C$は右辺blockを主要に解いた結果で，
$C(t_V)$は右辺targetへその代入を適用した型，$C(\Gamma)$はcontextへ同じ代入を適用したものを表す．
$\mathsf{gen}$はcontextに残っていない変数を量化してschemeを作る一般化である．
$s_b$は右辺終了supplyと閉じたcontextが要求する開始番号の大きい方である．
$\mathsf{interface}(\Gamma,C)$は，右辺を閉じた影響のうち外側contextから観測できる有限な等式列である．

\subsection{data constructor，primitive，条件式}

data constructorとprimitiveは，宣言表から閉じたschemeを引き，instantiateした関数型へ
引数を一つずつ\textsc{App}と同じ形で適用する．

\[
\frac{
  \Sigma(c)=\sigma\quad |\vec e|=\mathsf{arity}(\sigma)\quad
  \sigma\text{は閉じている}\quad
  \mathsf{Call}(\mathsf{instantiate}(\sigma,s),\vec e)=G;s'
}{
  \Sigma;\Gamma;s\vdash c(\vec e)\Longrightarrow G;s'
}\;\textsc{Ctor}
\]

\[
\frac{
  \Sigma(op)=\sigma\quad |\vec e|=\mathsf{arity}(\sigma)\quad
  \sigma\text{は閉じている}\quad
  \mathsf{Call}(\mathsf{instantiate}(\sigma,s),\vec e)=G;s'
}{
  \Sigma;\Gamma;s\vdash op(\vec e)\Longrightarrow G;s'
}\;\textsc{Prim}
\]

閉じたschemeとは，外から与えられていない自由変数を含まないschemeである．
\code{if c then t else e}は，組込みscheme
$\mathsf{Bool}\fn\alpha\fn\alpha\fn\alpha$をinstantiateし，$[c,t,e]$を同じ
$\mathsf{Call}$で処理する\textsc{If}規則である．これにより条件はBool，二枝は同じ型になる．

\section{M4のuser pattern規則}

user patternは，値に照合する\code{matchAll}や\code{matchFirst}のpatternである．
次の判断を使う．

\[
  \Sigma;\Gamma;\Phi;\Delta;s\vdash_p p
  \Longrightarrow\langle\delta;\Delta';H;P\rangle;s'
\]

$\Phi$はpattern functionへ渡されたpattern引数の型列，$\Delta$はそのpatternの左側ですでに
束縛された通常変数の型列である．$\delta=(\kappa\triangleright\tau)$はdualであり，
「capability $\kappa$を使って型$\tau$の値へ照合するpattern」を表す．$\Delta'$はこのpatternまでに
束縛された変数をsource順に並べた結果である．

\begin{longtable}{p{24mm}p{48mm}p{65mm}}
\toprule
pattern & 結果 & 説明 \\
\midrule
\endfirsthead
\toprule
pattern & 結果 & 説明 \\
\midrule
\endhead
variable & $\kappa_s\triangleright\alpha_s$，
$\Delta'=\Delta\mathbin{+\!+}[\alpha_s]$ & freshな対象型とcapabilityを作り，対象値を新しい変数として束縛する．\\
wildcard & $\kappa_s\triangleright\alpha_s$，$\Delta'=\Delta$ & 同じfresh dualを作るが，値を束縛しない．\\
value pattern $\#e$ & $\kappa_s\triangleright t_e$ & $e$を$\Delta\mathbin{+\!+}\Gamma$で調べ，
そのhardとpendingを引き継ぐ．新しいcapabilityだけを割り当てる．\\
constructor $c(p_1,\ldots,p_n)$ & 宣言済みdual schemeのresult & constructorをlookupし，arityを確認する．
各field patternのtargetとcapabilityを，instantiateしたfield dualへhard等式でそろえる．\\
tuple $(p_1,\ldots,p_n)$ &
$\Prod[\kappa_1,\ldots,\kappa_n]\triangleright
 \Prod[\tau_1,\ldots,\tau_n]$ & field patternを左から右へ処理し，capabilityとtargetをそれぞれproductにする．\\
conjunction $p\mathbin{\&}q$ & $\delta_p$と，$q$までのbinding列 & 同じ対象へ$p$，$q$の順で照合する．$p$のbindingを追加してから$q$を調べ，
両側のtarget型とcapabilityをhard等式で同じにする．\\
embedded parameter & $\Phi[i]$ & pattern functionの$i$番目の引数dualをそのまま使う．fresh変数や束縛を作らない．\\
pattern function $f(\vec p)$ & 宣言済みdual schemeのresult & frozen signatureから$f$のschemeだけをinstantiateする．
本体を再検査せず，field dualをschemeの引数dualへhard等式でそろえる．\\
\bottomrule
\end{longtable}

\section{M4で追加された式の規則}

M4の公開関係\code{M4.Elaborates}は，前節までの式規則を再帰的に使いながら，次の四形式を追加する．
内部では構文の大きさより大きい自然数をfuelとして持つ．fuelとは再帰呼出しの残り回数を表す数で，
公開判断は「十分なfuelが一つ存在する」とだけ要求するため，programの型には現れない．

\subsection{単項再帰\code{fix}}

\[
\frac{
  \mathsf{DirectSelf}(b)\qquad
  \Sigma;(d,\ d\fn r,\ \Gamma);s_b\vdash b
    \Longrightarrow\langle t_b;H_b;P_b\rangle;s'
}{
  \Sigma;\Gamma;s\vdash\mathsf{fix}\ b
  \Longrightarrow\langle d\fn r;H_b\mathbin{+\!+}[t_b=r];P_b\rangle;s'
}\;\textsc{Fix}
\]

body contextの先頭は引数$d$，次が再帰関数自身$d\fn r$である．\code{DirectSelf}は，
自己参照を別名に保存したり高階の値として逃がしたりせず，許可された直接呼出しだけを使う構文検査である．
bodyが構文上matcher literalなら，$d$をmatcher slot，$r$をmatcher producerの形へ最初から固定する．
これはmatcherという明示構文から決まる形であり，pendingから型の外形を推測する操作ではない．

\subsection{matcher literal}

\[
\frac{
  \mathsf{staticChecks}(\Sigma,\mathit{clauses})=\mathsf{true}
  \qquad
  \Sigma;\Gamma;\alpha_s;s_1\vdash_{clauses}
    \mathit{clauses}\Longrightarrow(E;H;P);s'
}{
  \Sigma;\Gamma;s\vdash\mathsf{matcher}\{\mathit{clauses}\}
  \Longrightarrow
  \langle\Matcher\,\kappa_s\,\alpha_s;
  \mathsf{evidenceEq}(\kappa_s,E)\mathbin{+\!+}H;P\rangle;s'
}\;\textsc{Matcher}
\]

$\alpha_s$はmatcherが調べる対象値の型，$\kappa_s$はmatcher全体のcapabilityである．
\code{staticChecks}は，constructorのarity，変数の重複禁止，catch-allの位置，data armの網羅性など，
型等式を解く前に判定できる構文条件をまとめた検査である．$E$は各clauseが示した分解能力の証拠で，
$\mathsf{evidenceEq}$がそれらを全体capabilityへhard等式として接続する．

一つのclauseは次の順で調べる．

\begin{enumerate}[leftmargin=2em]
  \item pattern-pattern headerを調べ，hole dual列とcapture型列を得る．holeは次に必要なmatcherの場所，
        captureは\code{\#\$x}のように対象値全体を式で参照する束縛である．
  \item next-matcher式を，holeが0個ならunit product，1個ならそのslot，2個以上ならslotのtupleとして
        checkingする．
  \item 各data-pattern armを対象型$\alpha_s$に対して調べ，arm bodyを
        $\mathsf{List}(\mathsf{holeProductTarget}(\mathit{holes}))$としてcheckingする．
\end{enumerate}

pattern-pattern headerのholeは，与えられた対象型とfreshなcapabilityを一組のhole dualにする．
constructor headerは宣言済みpattern-constructor schemeをinstantiateし，fieldを再帰的に調べる．
以下がheaderとarm patternのconstructor別一覧である．

\begin{longtable}{p{30mm}p{45mm}p{63mm}}
\toprule
種類 & constructor & 生成内容 \\
\midrule
\endfirsthead
\toprule
種類 & constructor & 生成内容 \\
\midrule
\endhead
pattern-pattern & hole & 与えられた対象型$\tau$とfreshな$\kappa_s$から
$[\kappa_s\triangleright\tau]$というhole列を作る．expected capabilityがあれば，
$\kappa_s$とのhard等式も追加する．\\
pattern-pattern & wildcard & holeもcaptureもhardも作らない．\\
pattern-pattern & capture & 与えられた対象型$\tau$をcapture列へ一つ追加する．holeは作らない．\\
pattern-pattern & constructor & 宣言済みschemeをinstantiateし，arityを確認してfieldを再帰的に調べる．
result targetを与えられた対象型へhard等式でそろえ，result capabilityを分解能力の証拠として返す．\\
data-pattern & variable & 与えられた対象型$\tau$をarm body用の束縛列へ追加する．\\
data-pattern & wildcard & 束縛もhardも作らない．\\
data-pattern & constructor & data-constructor schemeをinstantiateして関数型をfield型列とresult型へ分ける．
result型を与えられた対象型へhard等式でそろえ，fieldを再帰的に調べる．\\
data-pattern & tuple & 要素数だけfreshなfield型を作り，与えられた対象型をそのproductへhard等式でそろえる．
各要素を対応するfield型に対して再帰的に調べる．\\
\bottomrule
\end{longtable}

\subsection{全結果を返す\code{matchAll}}

target式$e_t$，matcher式$e_m$，user pattern $p$，body $e_b$の生成結果をそれぞれ
$T,M,Q,B$とする．$Q$のdualを$\kappa_p\triangleright\tau_p$とすると，

\[
\frac{
  \begin{gathered}
  \Sigma;\Gamma;s\vdash e_t\Longrightarrow T;s_1,\qquad
  \Sigma;\Gamma;[];[];s_1\vdash_p p\Longrightarrow Q;s_2,\\
  \Sigma;\Gamma;s_2\vdash e_m\Longrightarrow M;s_3,\qquad
  \Sigma;(\Delta_Q\mathbin{+\!+}\Gamma);s_3\vdash e_b\Longrightarrow B;s'
  \end{gathered}
}{
  \Sigma;\Gamma;s\vdash
  \mathsf{matchAll}\ e_t\ \mathsf{as}\ e_m\ \mathsf{with}\ p\to e_b
  \Longrightarrow G;s'
}\;\textsc{MatchAll}
\]

ここで$G$のtargetは$\mathsf{List}(t_B)$である．hardには
$\tau_p=t_T$を追加し，pendingには
$t_M\checks\Slot\,\kappa_p\,t_T$を追加する．つまりpatternの対象型は実際のtarget値と等しくし，
matcher式はそのpatternが要求するslotとして使えるかを後でcheckingする．$\Delta_Q$はpatternが
束縛した変数列で，body contextの先頭へsource順に追加する．

\subsection{最初の結果を返す\code{matchFirst}}

\[
\frac{
  \begin{gathered}
  \mathsf{armsExhaustive}(\mathit{arms})=\mathsf{true},\\
  \Sigma;\Gamma;s\vdash e_t\Longrightarrow T;s_1,\qquad
  \Sigma;\Gamma;s_1\vdash e_m\Longrightarrow M;s_2,\\
  \Sigma;\Gamma;t_T;t_M;s_2\vdash_{arms}\mathit{arms}
    \Longrightarrow A;s'
  \end{gathered}
}{
  \Sigma;\Gamma;s\vdash
  \mathsf{matchFirst}\ e_t\ \mathsf{as}\ e_m\ \mathsf{with}\ \mathit{arms}
  \Longrightarrow\mathsf{fromMatchFirst}(T,M,A);s'
}\;\textsc{MatchFirst}
\]

arm列は空であってはならず，最後のpatternは必ず一致すると構文的に分かる形でなければならない．
各armではpattern targetを$t_T$へhard等式でそろえ，matcher $t_M$をpattern capabilityのslotとして
pending checkingする．最初のarm bodyのtargetを全体結果型とし，後続bodyのtargetはその型へhard等式で
そろえる．\code{matchAll}と違い，結果へ\code{List}を付けない．

\section{最終的な\code{Typing}規則}

制約生成に成功しただけでは，hardとpendingが解けるとは限らない．現行の最終判断は次の二段階である．

\[
\frac{
  \begin{gathered}
  \Sigma;\Gamma;\mathsf{initialSupply}(\Gamma)\vdash e
    \Longrightarrow G;s',\\
  C\text{は}G\text{の吸収的な主要closure},\qquad
  \pi=C(t_G)
  \end{gathered}
}{
  \Sigma;\Gamma\vdash_{\mathrm{principal}}e:\pi
}\;\textsc{PrincipalTyping}
\]

\[
\frac{
  \Sigma;\Gamma\vdash_{\mathrm{principal}}e:\pi
  \qquad \tau=S(\pi)
}{
  \Sigma;\Gamma\vdash e:\tau
}\;\textsc{Typing}
\]

$S$は型変数への任意の代入である．したがって\code{Typing}は主要型$\pi$そのものだけでなく，
その型変数をさらに具体化して得られる型$\tau$も認める．M2--M3では\code{Source.Typing}，
M4では\code{Source.M4.Typing}がこの同じ二段階の形を持つ．

\begin{warningbox}{一覧の範囲}
ここに載せたのは現行source ASTの型付け・制約生成規則である．hardの単一化，pendingのpromotionと
residual checking，closureの主要性を証明する規則群は，生成されたblockを解く側の関係であり，
source構文の型規則とは分けている．またM4の健全性は実装済みだが，M4全体について
「宣言的に型が付けば公開推論が必ず成功する」という完全性と，すべての導出間の主要性は未完である．
この一覧は未証明の完全性を仮定しない．
\end{warningbox}

\chapter{用語集}

\begin{description}[leftmargin=10em,style=nextline]
  \item[absorbing closure]
    後続の解を合成しても，選んだ主要代表が正規化された形で保たれるclosure．
  \item[branch-stable acceptance]
    hardの主要解でchecking分岐を固定し，後から分岐選択をやり直さない受理概念．
  \item[capability]
    matcherが入力値のどの形を観察できるかを表す型情報．
  \item[closure]
    生成blockのhardと残余checkingを主要に解く代入と結果をまとめたもの．実行時関数closureとは別．
  \item[coherence]
    同じsourceに対する異なる正しい導出が，受理と主要結果について整合する性質．
  \item[entailed]
    あるhard理論を満たすすべての代入の下で必ず成り立つこと．
  \item[frame]
    生成済み制約blockを穴へ入れる周囲の制約構造．
  \item[freshness]
    新しく割り当てた変数名が，既存のsupportと衝突しない性質．
  \item[full-cut]
    \code{let}右辺の制約blockをbodyより先に完全に閉じる方式．
  \item[hard]
    checking変換の種類に依存せず必ず満たす型・capability等式．
  \item[interface equations]
    閉じた\code{let}右辺が外側contextの自由変数へ与えた影響を表す有限な等式列．
  \item[matcher producer]
    実際に作られたmatcher値の型．
  \item[matcher slot]
    matcherを使用する場所がmatcherへ課す要求型．
  \item[pending]
    source型をexpected型として使えるか，hardを解いた後で分類するchecking要求．
  \item[principal]
    ほかのすべての解が代入を通して得られるという意味で最も一般的であること．
  \item[replay]
    関係的導出から，公開推論器が実際に選ぶclosureを用いた導出を再構成すること．
  \item[scheme]
    量化された型変数を持つ多相型．現行M2では量化変数を位置で表す．
  \item[supply]
    新しい型変数名とcapability変数名を割り当てる二つの自然数カウンタ．
  \item[support]
    型，制約，代入などに実際に現れる有限な変数集合．
  \item[target]
    制約生成直後の生の結果型．最終代入の適用前なので型変数を含んでよい．
  \item[terminal audit]
    推論完了後に履歴や状態を再検査する旧方式．現行\code{Typing}の定義には含めない．
  \item[visible closure graph]
    二closureの有限な変数対応のうち，外側closed contextから観測できる移動辺だけを集めた証明用構造．
  \item[well-formed supply]
    contextに既に現れる全自由変数より後からfresh名を割り当てる開始条件．
\end{description}

\chapter{旧説明から現行説明への対応}

\begin{longtable}{p{48mm}p{92mm}}
\toprule
旧資料で見かける語 & 現行\code{type-pm-mech2}での扱い \\
\midrule
\endfirsthead
\toprule
旧資料で見かける語 & 現行\code{type-pm-mech2}での扱い \\
\midrule
\endhead
\code{SourceTyping} & 使用しない．現行は実行可能推論器と独立した\code{Source.Typing}を定義する．\\
\code{inferRaw} & 旧API．現行source側は\code{elaborate}，\code{infer}，M4側は\code{M4.infer}を使う．\\
\code{wBridgeCheck} & 旧terminal auditの一部．現行の宣言的型付け定義には持ち込まない．\\
\code{TypingInvariant} & 旧推論履歴の検査．現行M2完全性はentailed alignment，closure alignment，visible graphで証明する．\\
\code{PublicTheorems.lean} & 旧リポジトリの索引．現行の公開M2入口は\path{Source/FullM2Completion.lean}．\\
\bottomrule
\end{longtable}

\backmatter

\chapter*{版の基準}
\addcontentsline{toc}{chapter}{版の基準}

本書は2026年8月20日に，\code{type-pm-mech2}の現行\code{main}の
実装と公開監査項目を基準に更新した．
文書追加以外の実装変更が入った場合は，設計と進捗を統合した\code{README.md}，
\code{TypePM/AxiomAudit.lean}と照合して更新する．

\end{document}
